%&10pt
\font\ttsmall=omtt8

\def\setsmallprinting{\ttsmall \let\it=\itsmall \let\rm=\rmsmall 
   \def\ttstrut{\vrule height8pt depth3pt width0pt}%
   \offinterlineskip \parskip=-1pt\relax
}

\def\uncatcodespecials{\def\do##1{\catcode`##1=12}\dospecials}

\def\setupverbatim{\setsmallprinting
  \def\par{\leavevmode\endgraf}
  \catcode`\`=\active
  \catcode`\"=12 % unactivate
  \obeylines
  \uncatcodespecials
  \obeyspaces
  \everypar{\ttstrut}}
{\obeyspaces\global\let =\ } % let active space = control space
{\catcode`\`=\active \gdef`{\relax\lq}}

\def\listing#1{\begingroup\setupverbatim\input#1 \endgroup}

\hsize=210mm \advance\hsize by-3cm
\vsize=320mm
\pdfhorigin=1.5cm
\nopagenumbers

\centerline{\font\bf=omb10 at 17pt \bf Оповещение об авариях М-200}
\bigskip
\bigskip

\newcount\n
\def\N{\advance\n by1\indent\hbox to0pt{\hskip-\parindent\bf\the\n.\hfil}}

\N
При каждом возникновении аварийного события М-200 запускается звонок по телефону 42-08.
{\sl Asterisk\/} будет звонить бесконечно пока не поднимут трубку.
При поднятии трубки проговаривается сообщение об аварии.
\smallskip

\N
Первым запускается bat-файл, который указывается в настройках программы {\sl SMPAlarm}.
Этот bat-файл заносит время срабатывания и сообщение в лог-файл и вызывает {\tt alarm.sh} на
компьютере с {\sl asterisk},
под пользователем {\tt user}, по ключу. Ключ генерируется программой
{\sl genkey\/} из пакета PuTTY, затем добавляется
в файл {\tt\char`~/.ssh/authorized\char`_keys2} на компьютере с {\sl asterisk}.

При первоначальной настройке соединения требуется подтвердить ssh-сервер.
Для этого нужно в тревоге указать bat-файл в котором просто
запускается {\sl putty\/} и вызвать срабатывание этого bat-файла;
затем вручную соединиться с ssh-сервером и подтвердить, и выйти.
\smallskip

\N
Для каждой АТС свой bat-файл. Отличаются только
именем лог-файла и названиями АТС в имени файла блокировки (см.\ ниже) и ssh-команде.
Название АТС в ssh-команде используется в
имени звукового файла.
Вот пример для {\tt m200upats.bat} (ssh-команда должна идти в одну строку; также
обратить внимание на формат конца строк):
\medskip
\hrule
\smallskip
\listing{m200upats.bat}
\smallskip
\hrule
\bigskip

\N
Иногда соединение
начинает переходит из рабочего состояния в нерабочее и наоборот. До тех пор
пока проблема не будет устранена, бывает желательным заблокировать срабатывание обзвона.
Для этого просто нужно закрыть соответствующую программу {\sl SMPAlarm}.
\bigskip

\N
Открыть доступ по сети на сервере АТС. Добавить в настройки {\sl iptables\/} следующее:
\smallskip
\centerline{\tt -A INPUT -p tcp --dport 22 -s 192.168.40.141 -j ACCEPT}
\medskip

\N
Скрипт {\tt alarm.sh} заносит время срабатывания и сообщение в лог-файл и запускает обзвон.

Вот скрипт {\tt /etc/asterisk/alarm.sh}:
\medskip
\hrule
\smallskip
\listing{alarm.sh}
\smallskip
\hrule
\bigskip

\N
Логи будут в файле {\tt /var/log/alarm.log}. Для этого нужно создать файл
{\tt /etc/rsyslog.d/alarm.conf} со следующим содержимым:
\smallskip
\centerline{\tt local0.alert  /var/log/alarm.log}
\medskip

\N
Скрипт {\tt alarm.sh} запускается под пользователем {\tt user},
поэтому чтобы запустить call-файл под пользователем {\tt asterisk}\footnote*{В случае
если {\sl asterisk\/}
установлен не из дистрибутива (т.е.~работает под {\tt root}), `{\tt (asterisk)}' здесь
убрать и в {\tt alarm.sh} вместо
`{\tt sudo -u asterisk}' писать `{\tt sudo}'; а также убедиться что права на запись в файл
{\tt call.sh}
только у рута и что он лежит в директории с правами только рута.}
(под которым работает {\sl asterisk\/}),
нужно прописать через {\tt sudo visudo} следующее:
\smallskip
\centerline{\tt user ALL = (asterisk) NOPASSWD: /etc/asterisk/call.sh}
\medskip

\N
Вот скрипт {\tt /etc/asterisk/call.sh}:
\medskip
\hrule
\smallskip
\listing{call.sh}
\smallskip
\hrule

\bye
